\section{Conclusion: From Competition to Living Benchmark}
\label{sec:living-benchmark}

The \pokeagent Challenge transitions from the NeurIPS 2025 competition to a living benchmark available at \url{https://pokeagentchallenge.com}. The Battling Track maintains a live leaderboard on a dedicated Showdown server with all organizer baselines active, allowing new agents to be evaluated against the full history of submissions. The Speedrunning Track provides self-contained evaluation: researchers run agents locally against the standardized emulator and milestone framework, enabling reproducible comparison without server access. All datasets, baselines, and infrastructure are publicly released. The evaluation server and leaderboards are maintained by the organizing team with funding support as listed in the Acknowledgements; all code and data are hosted on GitHub and HuggingFace to ensure long-term availability independent of server infrastructure. Both tracks provide clear room to expand in order to maintain benchmark difficulty as performance saturates.

Large capability gaps remain. We highlight four open challenges: \textbf{(1) VLM-SLAM:} Speedrunning agents struggle with basic localization, action-distance estimation, and objective detection. Grounding VLM outputs to consistent spatial representations---analogous to classical SLAM but through language-vision interfaces---remains a bottleneck for RPG play. \textbf{(2) Closing the LLM--RL gap in battling:} Specialist RL agents far outpace harness LLM agents in competitive battles. Developing LLM agents that match RL performance, or hybrid approaches that combine RL's sample efficiency with LLM world knowledge, is an open problem. \textbf{(3) Full-game completion with open-source models:} Proprietary frontier models have completed Pokémon RPGs with heavy harness support, but no open-source model has done so. Achieving this would make long-horizon RPG evaluation accessible to more research groups. \textbf{(4) Approaching human speedrun times:} The best agent (Heatz, 40:13) is $2.2\times$ slower than human speedrunners. Closing this gap requires advances in navigation efficiency, obstacle avoidance, and objective sequencing---capabilities relevant to time-critical planning more broadly.
